\begin{abstract}
本文研究统一基地约束下的多无人机协同巡检路径优化问题。题目给出四个规模不同的巡检案例，I、II、III 级巡检点分别需要 3、2、1 次有效到达；无人机从坐标 $(0,0)$ 的基地于 8:00 同时出发，以 $55\ \mathrm{km/h}$ 巡航，并在完成任务后返回基地。本文将重复巡检任务展开为带唯一标识的任务集合，特别规定同一路线中连续相同物理点只计一次有效到达，因此 A-A-A 计 1 次而 A-B-A 计 2 次。

问题一建立带最大工作时间约束的 min--max 多旅行商模型。我们先由不可压缩服务时间和最远点往返时间给出理论下界，再从下界起逐一搜索首个通过独立验证的可行机队，并在固定机队内最小化最大工作时间。问题二继承问题一的机队和全部硬约束，采用字典序目标 $\operatorname{lexmin}(C,\delta,D)$，通过跨路线 relocate、swap、2-opt* 和 cross-exchange 改善工作时间均衡。问题三将附件中的圆形禁飞区扩展为时间依赖边函数，同时检查线段几何相交、飞行时间重叠、巡检服务时段和基地停留，并在固定机队不可行时逐架增机。

四个 Case 的问题一首个已验证可行机队数分别为 5、2、6、5 架，最大工作时间分别为 8.7129、7.9583、8.6125、8.1378 h；问题二保持上述机队，最大工作时间为 8.6139、7.8011、8.5980、8.1378 h，工作时间极差为 0.0896、0.0027、0.0976、0.0644 h；问题三的最小已验证可行机队数为 5、3、9、7 架，最大工作时间为 8.7400、7.1608、7.9812、8.3802 h，四个正式结果均通过任务覆盖、回基地、9 h 和动态禁飞独立校验。本文同时保留理论下界与启发式首个可行数量的区别，不将未获得精确证书的结果表述为全局最优。

\noindent\textbf{关键词：}多无人机协同巡检；有效到达；min--max 车辆路径；负载均衡；动态禁飞区；时空路径规划
\end{abstract}
